%% titre: Quiz 1 — Notions de logique
%% niveau: Licence 1
%% matiere: Algèbre 1
%% annee: 2026-2027
%% ordre: 1
%% description: Négations, implication, quantificateurs et modes de raisonnement.

\begin{question}
Quelle est la négation de $\forall x\in\R,\ \exists y\in\R,\ y>x$ ?
\begin{reponses}
  \mauvaise $\forall x\in\R,\ \exists y\in\R,\ y\le x$
  \bonne $\exists x\in\R,\ \forall y\in\R,\ y\le x$
  \mauvaise $\exists x\in\R,\ \exists y\in\R,\ y\le x$
  \mauvaise $\exists x\in\R,\ \forall y\in\R,\ y<x$
\end{reponses}
\begin{explication}
On échange chaque $\forall$ en $\exists$ (et inversement) en gardant l'ordre, puis on nie le prédicat : la négation de $y>x$ est $y\le x$.
\end{explication}
\end{question}

\begin{question}
Quelles propositions sont équivalentes à $P\Rightarrow Q$ ? (plusieurs réponses possibles)
\begin{reponses}
  \bonne $\overline{P}\vee Q$
  \bonne $\overline{Q}\Rightarrow\overline{P}$
  \mauvaise $Q\Rightarrow P$
  \mauvaise $P\wedge\overline{Q}$
\end{reponses}
\begin{explication}
$P\Rightarrow Q$ équivaut à $\overline P\vee Q$ et à sa contraposée $\overline Q\Rightarrow\overline P$. La proposition $Q\Rightarrow P$ est la réciproque, et $P\wedge\overline Q$ est la \emph{négation} de $P\Rightarrow Q$.
\end{explication}
\end{question}

\begin{question}
Si $P$ est fausse, la proposition $P\Rightarrow Q$ est :
\begin{reponses}
  \mauvaise fausse
  \bonne vraie, quelle que soit la valeur de vérité de $Q$
  \mauvaise vraie seulement si $Q$ est vraie
  \mauvaise sans valeur de vérité
\end{reponses}
\begin{explication}
$P\Rightarrow Q$ n'est fausse que dans le cas $P$ vraie et $Q$ fausse.
\end{explication}
\end{question}

\begin{question}
La négation de $P\wedge Q$ est :
\begin{reponses}
  \bonne $\overline{P}\vee\overline{Q}$
  \mauvaise $\overline{P}\wedge\overline{Q}$
  \mauvaise $P\vee Q$
  \mauvaise $\overline{P}\Rightarrow\overline{Q}$
\end{reponses}
\begin{explication}
C'est une loi de De Morgan : on nie chaque proposition et on échange $\wedge$ et $\vee$.
\end{explication}
\end{question}

\begin{question}
Laquelle de ces propositions est vraie ?
\begin{reponses}
  \mauvaise $\forall x\in\R,\ x^2>0$
  \mauvaise $\exists x\in\R,\ x^2=-1$
  \bonne $\forall x\in\R,\ \exists y\in\R,\ y>x$
  \mauvaise $\exists y\in\R,\ \forall x\in\R,\ y>x$
\end{reponses}
\begin{explication}
Pour tout réel $x$, le réel $y=x+1$ convient. En revanche aucun réel n'est strictement supérieur à tous les réels (y compris lui-même), et $0^2=0$ n'est pas $>0$.
\end{explication}
\end{question}

\begin{question}
Pour démontrer une proposition $P$ par l'absurde, on :
\begin{reponses}
  \bonne suppose que $P$ est fausse et on aboutit à une contradiction
  \mauvaise suppose que $P$ est vraie et on vérifie quelques exemples
  \mauvaise démontre la réciproque de $P$
  \mauvaise démontre $\overline{Q}\Rightarrow\overline{P}$
\end{reponses}
\begin{explication}
La dernière proposition décrit le raisonnement par contraposée, qui est différent.
\end{explication}
\end{question}

\begin{question}
Dans une démonstration par récurrence, l'étape d'hérédité consiste à :
\begin{reponses}
  \mauvaise vérifier que $\mathcal{P}(n_0)$ est vraie
  \bonne montrer que, pour tout $n\ge n_0$, $\mathcal{P}(n)\Rightarrow\mathcal{P}(n+1)$
  \mauvaise supposer $\mathcal{P}(n)$ vraie pour \emph{tout} entier $n$
  \mauvaise démontrer $\mathcal{P}(n+1)$ sans utiliser $\mathcal{P}(n)$
\end{reponses}
\begin{explication}
La première réponse décrit l'initialisation. On suppose $\mathcal P(n)$ pour un entier $n$ \emph{fixé}, pas pour tous.
\end{explication}
\end{question}
